\SourceLeaf{132}
\begin{quote}\textit{die unklare Fassung desjenigen Theils des Regulativs, wodurch die Progression der Pensionssätze bestimmt werden soll}\end{quote}
in Verbindung mit
\begin{quote}\raggedright\textit{der gegenwärtig so sehr vergrösserten Anzahl der an der Witwenkasse Theil nehmen en Professoren.}\end{quote}
\noindent\hspace*{\parindent}Die grosse Wichtigkeit des letztern Umstandes ist schon in meinem oben erwähnten Votum angedeutet. Die Bedeutsamkeit einer grossen Anzahl von Interessenten ist nach der Beschaffenheit einer Witwenkasse eine sehr verschiedene. Für eine Witwenkasse, welche sich durch die Beiträge der Mitglieder (oder durch die Antrittsgelder, oder durch beides verbunden) ganz selbst erhält, wird eine recht grosse Anzahl der Theilnehmer nur vortheilhaft sein, vorausgesetzt, dass die Kasse auf eine richtige Rechnung basirt ist. Eine doppelt starke Gesellschaft \textit{dieser} Art, hat unter übrigens gleichen Umständen eine doppelt so grosse Anzahl von Witwen zu erwarten, wie eine einfache: sie hat aber auch gerade doppelt so viele Einnahme, und kann daher den einzelnen Witwen gerade eben so viel gewähren, aber mit mehr Sicherheit gegen die wechselnden Fluctuationen, welche bei der grössern Gesellschaft im Verhältniss zum Ganzen geringer sind, als bei der kleinern.

Ganz anders aber verhält es sich mit einer Witwenverpflegungsanstalt, die ein reines Beneficium ist, und deren Mittel einmal eine gegebene Grösse haben (durch bestimmten Kapital- oder Grundbesitz). Auch hier wird jede Erweiterung des Umfanges eine in gleichem Verhältnisse vermehrte Anzahl der Witwen zur Folge haben, deren jede einzelne demnach auch nicht mehr so viel aus den gegebenen Mitteln wird erhalten können, wie vorher bei beschränktem Umfange. Allerdings wird die der vergrösserten Interessentenzahl entsprechende Vergrösserung der Witwenzahl in ihrer \textit{vollen} Stärke erst nach mehrern Decennien eintreten, und dem natürlichen Gange der Dinge gemäss bis dahin sich nach und nach entwickeln. Setzen wir, um die Vorstellung zu fixiren, der Umfang einer solchen Gesellschaft (die ein reines Beneficium ist) habe sich binnen einer gewissen Zeit verdoppelt. Man wird dann bald auf eine vergrösserte Witwenzahl, also, wenn das Vermögen nicht selbst angegriffen werden soll, auf eine Verminderung der jeder einzelnen Witwe zu gewährenden Pension gefasst sein müssen, und diese Herabsetzung wird nach und nach bis auf die Hälfte fortschreiten. Hätte aber eine solche Gesellschaft ein Statut, wonach den Witwen, trotz ihrer steigenden Zahl, fortwährend gleichbleibende Pensionen gezahlt werden \textit{müssen}, so würde sie nothwendig zu Grunde gehen. Zwei Fälle gibt es jedoch, wo dieser Hergang eine Modification erleiden wird oder erleiden kann. \textit{Erstlich} wenn die Mittel der Kasse, vor der Erweiterung des Umfanges der Theilnahme, mehr als hinreichend waren, um die bestehenden Pensionen zu bestreiten, so dass eine fortwährende Vermögensvergrösserung, und etwa auch bis dahin von Zeit zu Zeit eine Erhöhung des Pensionssatzes hatte Statt finden können. Hier wird offenbar der Erfolg des erweiterten Umfanges von dem \textit{Wieviel?} abhängen. Hatte die Kasse vorher einen grossen jährlichen Ueberschuss, und ist die Vergrösserung der Interessentenzahl nicht sehr bedeutend, so kann jene die Gefahr vielleicht überstehen; die Vermögenszunahme wird nur immer langsamer und langsamer werden, und möglicherweise kann, wenn die Folge jener Ursache sich erst ganz entwickelt hat, die Kraft der Kasse noch hinreichend sein, auch der grössern Witwenzahl die volle Pension zu gewähren. Umgekehrt aber, war anfänglich der jährliche Ueberschuss nicht gross, die Vermehrung der Interessentenzahl aber sehr erheblich, so wird der jährliche Ueberschuss bald in ein Deficit übergehen, und der Ruin der Kasse zwar etwas später, als wenn ursprünglich Mittel und Ansprüche im Gleichgewicht waren, aber doch eben so unfehlbar eintreten. \textit{Zweitens} bei einer Kasse von überhaupt geringem Umfange in Beziehung auf die Zahl der Theilnehmer, und wo diese Zahl also auch nach der Vergrösserung noch wie eine kleine zu betrachten ist, wird man keinen so regelmässigen Hergang erwarten dürfen wie bei grössern; die in der Natur der Sache liegenden, aber, bei kleinen Zahlen \textit{verhältnissmässig} viel grössern Schwankungen werden die Regelmässigkeit in der Folge der Erscheinungen
